% GENERATED FILE. Edit canonical lessons, then run npm run generate:books.
% canonical-source-hash: fnv1a64:4415afda2f664c6f
% canonical-lessons: JA-C211-doubutsu, JA-C211-saru, JA-C211-usagi, JA-C211-kotori, JA-C211-fueru

\chapter{Animals in the Hills}
\label{ch:ja-animals-in-the-hills}
\hlchaptermodality{211}

\begin{chapteropening}
\textbf{By the end of this chapter:} \emph{I can name some animals and say where they live.}

Complete the last lesson of chapter 211: I can name some animals and say where they live.
\end{chapteropening}

\section[\texorpdfstring{\ja{どうぶつ} (dōbutsu)}{どうぶつ (dōbutsu)}]{\ja{どうぶつ} (dōbutsu) — an animal}

\label{lesson:JA-C211-doubutsu}



\begin{quote}
Before the new one: say the Japanese for to grow, then the Japanese for to wither.
\end{quote}

\subsection*{You'll want to know: \ja{どうぶつ}}
\textbf{\ja{どうぶつ}} — \emph{dōbutsu} — \textquotedblleft{}an animal\textquotedblright{}.

\begin{quote}
\textbf{\ja{もりに} \ja{どうぶつが} \ja{います。}} — \emph{mori ni dōbutsu ga imasu} — \textquotedblleft{}There are animals in the forest.\textquotedblright{}
\end{quote}

\begin{cousinweb}
\textbf{\ja{どうぶつ}} (\emph{dōbutsu}) is Sino-Japanese (kango): two Chinese characters read with their on'yomi, readings borrowed from Chinese. They are \emph{dō}, \textquotedblleft{}to move\textquotedblright{} (Middle Chinese \emph{duwngX}) and \emph{butsu}, \textquotedblleft{}thing\textquotedblright{} (\emph{mjut}) (KANJIDIC2; Middle Chinese in Baxter's transcription).
\end{cousinweb}

\begin{grammarlens}[title={What you've built}]
One more word for everyday life at A2.
\end{grammarlens}

\subsection*{Your turn}
\begin{itemize}
\raggedright
  \item \emph{Say it:} \emph{dōbutsu}
  \item \emph{Say it:} \emph{dōbutsu}, once more
  \item \emph{Say it:} \emph{mori ni dōbutsu ga imasu} — There are animals in the forest.
  \item \emph{Recall:} say the Japanese for to grow, then the Japanese for to wither, then say \emph{dōbutsu} again
\end{itemize}

\begin{tcolorbox}[breakable,colback=teal!4,colframe=teal!35!black,title={Before you move on}]
What does \ja{どうぶつ} mean? (\textquotedblleft{}An animal\textquotedblright{}.) Say it once more.
\end{tcolorbox}
\section[\texorpdfstring{\ja{さる} (saru)}{さる (saru)}]{\ja{さる} (saru) — a monkey}

\label{lesson:JA-C211-saru}



\begin{quote}
Before the new one: say the Japanese for to wither, then the Japanese for an animal.
\end{quote}

\subsection*{You'll want to know: \ja{さる}}
\textbf{\ja{さる}} — \emph{saru} — \textquotedblleft{}a monkey\textquotedblright{}.

\begin{quote}
\textbf{\ja{やまに} \ja{さるが} \ja{います。}} — \emph{yama ni saru ga imasu} — \textquotedblleft{}There are monkeys on the mountain.\textquotedblright{}
\end{quote}

\begin{cousinweb}
\textbf{\ja{さる}} (\emph{saru}) is a native Japanese word (wago), inherited rather than borrowed from Chinese; its kanji was attached later as a kun'yomi, and its older origin is not settled (JMdict, KANJIDIC2).
\end{cousinweb}

\begin{grammarlens}[title={What you've built}]
One more word for everyday life at A2.
\end{grammarlens}

\subsection*{Your turn}
\begin{itemize}
\raggedright
  \item \emph{Say it:} \emph{saru}
  \item \emph{Say it:} \emph{saru}, once more
  \item \emph{Say it:} \emph{yama ni saru ga imasu} — There are monkeys on the mountain.
  \item \emph{Recall:} say the Japanese for to wither, then the Japanese for an animal, then say \emph{saru} again
\end{itemize}

\begin{tcolorbox}[breakable,colback=teal!4,colframe=teal!35!black,title={Before you move on}]
What does \ja{さる} mean? (\textquotedblleft{}A monkey\textquotedblright{}.) Say it once more.
\end{tcolorbox}
\section[\texorpdfstring{\ja{うさぎ} (usagi)}{うさぎ (usagi)}]{\ja{うさぎ} (usagi) — a rabbit}

\label{lesson:JA-C211-usagi}



\begin{quote}
Before the new one: say the Japanese for an animal, then the Japanese for a monkey.
\end{quote}

\subsection*{You'll want to know: \ja{うさぎ}}
\textbf{\ja{うさぎ}} — \emph{usagi} — \textquotedblleft{}a rabbit\textquotedblright{}.

\begin{quote}
\textbf{\ja{うさぎは} \ja{くさを} \ja{たべます。}} — \emph{usagi wa kusa o tabemasu} — \textquotedblleft{}Rabbits eat grass.\textquotedblright{}
\end{quote}

\begin{cousinweb}
\textbf{\ja{うさぎ}} (\emph{usagi}) is a native Japanese word (wago), inherited rather than borrowed and usually written in kana; its kanji was attached later as a kun'yomi, and its older origin is not settled (JMdict, KANJIDIC2).
\end{cousinweb}

\begin{grammarlens}[title={What you've built}]
One more word for everyday life at A2.
\end{grammarlens}

\subsection*{Your turn}
\begin{itemize}
\raggedright
  \item \emph{Say it:} \emph{usagi}
  \item \emph{Say it:} \emph{usagi}, once more
  \item \emph{Say it:} \emph{usagi wa kusa o tabemasu} — Rabbits eat grass.
  \item \emph{Recall:} say the Japanese for an animal, then the Japanese for a monkey, then say \emph{usagi} again
\end{itemize}

\begin{tcolorbox}[breakable,colback=teal!4,colframe=teal!35!black,title={Before you move on}]
What does \ja{うさぎ} mean? (\textquotedblleft{}A rabbit\textquotedblright{}.) Say it once more.
\end{tcolorbox}
\section[\texorpdfstring{\ja{ことり} (kotori)}{ことり (kotori)}]{\ja{ことり} (kotori) — a small bird}

\label{lesson:JA-C211-kotori}



\begin{quote}
Before the new one: say the Japanese for a monkey, then the Japanese for a rabbit.
\end{quote}

\subsection*{You'll want to know: \ja{ことり}}
\textbf{\ja{ことり}} — \emph{kotori} — \textquotedblleft{}a small bird\textquotedblright{}.

\begin{quote}
\textbf{\ja{あさ} \ja{ことりが} \ja{なきます。}} — \emph{asa kotori ga nakimasu} — \textquotedblleft{}Small birds sing in the morning.\textquotedblright{}
\end{quote}

\begin{cousinweb}
\textbf{\ja{ことり}} (\emph{kotori}) is native Japanese (wago), built from \emph{ko}, \textquotedblleft{}small\textquotedblright{}, and \emph{tori}, \textquotedblleft{}bird\textquotedblright{}, both inherited rather than borrowed (JMdict, KANJIDIC2).
\end{cousinweb}

\begin{grammarlens}[title={What you've built}]
One more word for everyday life at A2.
\end{grammarlens}

\subsection*{Your turn}
\begin{itemize}
\raggedright
  \item \emph{Say it:} \emph{kotori}
  \item \emph{Say it:} \emph{kotori}, once more
  \item \emph{Say it:} \emph{asa kotori ga nakimasu} — Small birds sing in the morning.
  \item \emph{Recall:} say the Japanese for a monkey, then the Japanese for a rabbit, then say \emph{kotori} again
\end{itemize}

\begin{tcolorbox}[breakable,colback=teal!4,colframe=teal!35!black,title={Before you move on}]
What does \ja{ことり} mean? (\textquotedblleft{}A small bird\textquotedblright{}.) Say it once more.
\end{tcolorbox}
\section[\texorpdfstring{\ja{ふえる} (fueru)}{ふえる (fueru)}]{\ja{ふえる} (fueru) — to increase; to multiply}

\label{lesson:JA-C211-fueru}



\begin{quote}
Before the new one: say the Japanese for a rabbit, then the Japanese for a small bird.
\end{quote}

\subsection*{You'll want to know: \ja{ふえる}}
\textbf{\ja{ふえる}} — \emph{fueru} — \textquotedblleft{}to increase; to multiply\textquotedblright{}.

\begin{quote}
\textbf{\ja{らいねん} \ja{うさぎが} \ja{ふえます。}} — \emph{rainen usagi ga fuemasu} — \textquotedblleft{}Next year there will be more rabbits.\textquotedblright{}
\end{quote}

Polite \textbf{\ja{ふえます}}, \emph{fuemasu}; an ichidan verb.

\begin{cousinweb}
\textbf{\ja{ふえる}} (\emph{fueru}) is a native Japanese word (wago), inherited rather than borrowed from Chinese; its kanji was attached later as a kun'yomi, and its older origin is not settled (JMdict, KANJIDIC2).
\end{cousinweb}

\begin{grammarlens}[title={What you've built}]
One more word for everyday life at A2.
\end{grammarlens}

\subsection*{Your turn}
\begin{itemize}
\raggedright
  \item \emph{Say it:} \emph{fueru}
  \item \emph{Say it:} \emph{fueru}, once more
  \item \emph{Say it:} \emph{rainen usagi ga fuemasu} — Next year there will be more rabbits.
  \item \emph{Recall:} say the Japanese for a rabbit, then the Japanese for a small bird, then say \emph{fueru} again
\end{itemize}

\begin{tcolorbox}[breakable,colback=teal!4,colframe=teal!35!black,title={Before you move on}]
What does \ja{ふえる} mean? (\textquotedblleft{}To increase; to multiply\textquotedblright{}.) Say it once more.
\end{tcolorbox}
